% §17.4 Sidebar: Fully Homomorphic Encryption
%
% Cross-reference resolution status:
%
% DEFINED here:
%   \label{sec:fhe-sidebar}                   -- this section
%
% ALREADY DEFINED (chapter 5):
%   \Cref{sec:bloom-from-scratch}             -- chapter 5 §5.2 Bloom filter
%
% ALREADY DEFINED (chapter 14):
%   \Cref{ch:optimal-bernoulli-sets}          -- chapter 14 BHF construction
%
% INTERNAL refs (chapter 17):
%   \Cref{def:cipher-map}                     -- §17.1
%   \Cref{sec:trapdoor-paradigm}              -- §17.3
%
% §17.4 is the FHE sidebar. Five items: (1) what FHE is, Gentry
% (2009); (2) cipher maps as a weaker, fprate/fnrate abstraction;
% (3) FHE cost (huge ciphertexts, slow operations); (4) when to use
% each (FHE for absolute correctness; cipher maps for tolerable
% error plus performance); (5) concrete comparison (Bloom cipher map
% at epsilon=0.01 vs FHE equivalent, ~10^4 space factor).

\section{Sidebar: Fully Homomorphic Encryption}
\label{sec:fhe-sidebar}

The previous sections built up cipher maps as approximate functions
on encrypted data. The approximate qualifier is a deliberate
weakening: a cipher map accepts bounded $\fprate$ and $\fnrate$ as
the price for being practical. A natural question is what the
unweakened version looks like. The answer is fully homomorphic
encryption, the cryptographic abstraction that promises zero-error
computation on ciphertexts. This sidebar steps off the main thread
of the chapter to lay out what FHE is, where it sits relative to
cipher maps, and when each abstraction is the right tool for the
job.

\subsection*{What fully homomorphic encryption is}

Fully homomorphic encryption (FHE) is an encryption scheme in which
the ciphertexts support exact arithmetic. Given $\mathsf{Enc}(x)$
and $\mathsf{Enc}(y)$, an FHE scheme provides operations
$\mathsf{Enc}(x) \oplus \mathsf{Enc}(y) = \mathsf{Enc}(x + y)$ and
$\mathsf{Enc}(x) \otimes \mathsf{Enc}(y) = \mathsf{Enc}(x \cdot y)$,
preserved exactly through any depth of composition. Because
arbitrary Boolean and arithmetic circuits decompose into additions
and multiplications, FHE supports the exact evaluation of any
computable function on encrypted inputs: a server can compute $\tilde
f(\mathsf{Enc}(x_1), \ldots, \mathsf{Enc}(x_n))$ and return
$\mathsf{Enc}(f(x_1, \ldots, x_n))$ with zero error in the
plaintext output.

Gentry's 2009 thesis gave the first viable construction of an FHE
scheme, resolving a question that had stood open since Rivest,
Adleman, and Dertouzos posed it in 1978. Earlier schemes had been
\emph{partially} homomorphic, supporting either addition (Paillier)
or multiplication (RSA, ElGamal) but not both, and so could not
evaluate arbitrary circuits. Gentry's construction allowed both
operations and showed how to control the ciphertext noise that
accumulates with depth. Subsequent schemes (BFV, BGV, CKKS, TFHE)
have refined the construction, traded off different cost dimensions,
and made FHE practical for a growing class of applications.

For the present chapter the relevant fact is that FHE realizes a
cipher map with $\fprate = \fnrate = 0$. The client recovers the
exact plaintext output of $f$, with no error from the cipher
construction itself. Whatever uncertainty the client has about the
plaintext after decrypting an FHE result is uncertainty that was
already present in the inputs, not uncertainty introduced by the
computation.

\subsection*{Cipher maps as a weaker abstraction}

A cipher map in the sense of \Cref{def:cipher-map} is the
strictly weaker abstraction. The cipher map admits a positive
$\fprate$ and $\fnrate$, and the framework's job is to bound the
error and propagate it through compositions. What the cipher map
buys in exchange for accepting error is room to operate. The
designer of a cipher map can choose representations (hash families,
sketch data structures, Bloom filters) that are orders of magnitude
smaller and faster than the polynomial-arithmetic representations
FHE requires. The Bernoulli framework's space-accuracy duality,
which has been the backbone of chapters 5 through 15, applies
directly: a cipher map for set membership at $\fprate = 2^{-20}$
uses about $20 + \mu$ bits per element by the optimal construction
of \Cref{ch:optimal-bernoulli-sets}, while the corresponding FHE
realization (an exact encrypted set) uses kilobytes per element.

The two abstractions therefore occupy different rungs on the same
ladder. FHE is the strongest abstraction the cryptographer can
deliver and the most expensive to use. Cipher maps are weaker (they
allow error) but cheaper in proportion to how much error the
application tolerates. The framework's role is to make the trade-off
quantitative: a cipher map with $\fprate = \epsilon$ uses about
$-\log_2 \epsilon$ bits per element of useful information capacity,
the cipher-map equivalent of FHE's exact-but-expensive guarantee.

\subsection*{The cost of FHE}

The price of FHE's zero-error guarantee is steep in three dimensions.

\textbf{Ciphertext size.} An FHE ciphertext encrypts a single
plaintext bit but occupies many kilobytes of memory. Early schemes
ran in the megabyte-per-bit range; modern BFV, BGV, and CKKS
ciphertexts pack tens to hundreds of plaintext bits into one
ciphertext, but still cost on the order of kilobytes per encrypted
bit at typical security parameters (128 to 256 bits of security
against quantum or classical attack). For a database that would
fit comfortably in a few gigabytes of plaintext, the FHE
representation runs to terabytes.

\textbf{Operation latency.} Each homomorphic addition or
multiplication takes milliseconds on a modern processor; deeper
circuits compound the cost. A simple lookup that takes nanoseconds
on plaintext data can take seconds or minutes under FHE, even on
the fastest available implementations. The slowdown is roughly six
to nine orders of magnitude relative to the plaintext operation,
depending on the scheme and the circuit depth.

\textbf{Bootstrapping.} The ciphertext noise that accumulates with
each homomorphic operation must periodically be reset; the
\emph{bootstrapping} operation that performs this reset is by far
the most expensive primitive in FHE, taking seconds even in the
fastest current implementations. Practical FHE workloads either
budget bootstrapping carefully (running only depth-bounded circuits
between resets) or accept the bootstrapping cost as the dominant
runtime component.

The combined cost is why FHE has not displaced plaintext computation
even for security-critical applications. The cryptographic community
has made enormous progress since 2009 (the cost has dropped by many
orders of magnitude), and FHE is now practical for narrow workloads
(simple aggregations, fixed-circuit medical queries, contract
evaluations on small data). But for the broad class of workloads
that cipher maps address (approximate set membership, ranked
retrieval, anonymity-preserving analytics over large datasets), FHE
remains too slow and too large.

\subsection*{When to use which}

The choice between FHE and a cipher map is the choice between exact
and approximate. Two questions decide it.

\textit{Does the application tolerate error?} A contract evaluation
on encrypted inputs is exact by definition: a single bit-flip in
the output changes the meaning of the contract. A medical record
lookup that returns whether a patient has a specific diagnosis is
exact in the sense that a false positive could trigger a wrongful
treatment. FHE is the right choice for these correctness-critical
queries; cipher maps are not. A search query that returns the top
ten matching documents from a million-document corpus tolerates a
small false-positive rate (a few irrelevant documents in the top
ten are an inconvenience, not a catastrophe). A trapdoor cipher
map at $\fprate = 0.01$ is the right choice here.

\textit{Does the application need to scale?} A handful of exact
queries per day on a small encrypted dataset can run under FHE
without difficulty: even at minutes per query, the throughput is
adequate. A search index serving thousands of queries per second
on a multi-gigabyte dataset cannot run under FHE at any realistic
hardware budget. Cipher maps scale because they live in the same
asymptotic regime as plaintext data structures: a Bloom-filter
cipher map serves queries in constant time and uses space
proportional to the dataset size, the same complexity as a
cleartext Bloom filter.

In practice the two abstractions often combine. A hybrid system
might use a cipher map as the first-stage filter (rejecting most
non-matches quickly with bounded error) and FHE as the second-stage
verifier (exactly checking the cipher map's positive results). The
cipher map keeps the per-query cost low; the FHE stage cleans up
the false positives. The framework's composition theorem of
\Cref{thm:cipher-composition} applies to this hybrid: the
overall $(\fprate, \fnrate)$ is the matrix product of the two
stages, and the FHE stage's exact channel (identity matrix) drops
the false-positive rate to whatever the cipher map can verify.

\subsection*{A concrete comparison}

A worked example fixes the cost ratio. Take a set membership query
on a plaintext set $A$ of $|A| = 10^6$ elements, each represented
by a 64-bit integer.

\textbf{Bloom-filter cipher map at $\fprate = 0.01$.} The optimal
Bloom-filter construction of chapter 5 uses about $-\log_2(0.01)
\approx 6.64$ bits per element on the underlying filter, plus a
small constant overhead for hash-family parameters. The trapdoor
adds a keyed seed (a few dozen bytes total) but does not change the
per-element cost asymptotically. The total cipher-map representation
is about $10^6 \times 6.64 \approx 6.64 \times 10^6$ bits, or
roughly 830 kilobytes. Each membership query runs in constant time
(a few hash computations and bit lookups), with no bootstrapping or
ciphertext expansion overhead.

\textbf{FHE-encrypted exact set.} An FHE realization of the same
set membership query encrypts each element under a scheme like BFV
or BGV. At a security parameter of 128 bits, each ciphertext
occupies about $10^4$ to $10^5$ bits, or roughly 1 to 10 kilobytes
per encrypted element. The total representation is about $10^6
\times 10^4 = 10^{10}$ bits, or roughly 1 to 10 gigabytes for the
same 1 million elements. Each membership query under FHE requires
homomorphic comparisons against each ciphertext or against a
homomorphic hash, taking seconds in either case.

\textbf{The ratio.} The cipher map costs about 830 kilobytes; the
FHE version costs about 1 to 10 gigabytes. The space ratio is on
the order of $10^4$, with the cipher map four orders of magnitude
smaller. The runtime ratio is roughly the same: nanoseconds for the
cipher map, seconds for FHE, giving about $10^9$ in favor of the
cipher map. The four-orders-of-magnitude space gap and nine-orders
runtime gap are the price the cipher map pays in correctness (a
$\fprate$ of 1 percent) for performance practical enough to serve a
production search engine.

This is the central trade-off the trapdoor computing paradigm
exploits. FHE remains the abstraction of choice when error is
intolerable. Cipher maps fill the much larger space of workloads
where bounded error is acceptable, and the framework's job is to
keep the error rates honest as the cipher maps compose. The next
section returns to the cipher-map setting and asks how Boolean
combinations of cipher maps behave under the framework's
composition rules.
